% TrArrive/TrAwait are paired side by side to make the arrive/await pairing clear.
% Fence/Req diagrams put their quadrants on the sides to save vertical space.
\begin{figure}[h]
\sporkCodeRule
\centering
\inOutDiagram
  {\sporkTrPostAccess(\sporkmut, \sporkx, \sporkz_i^*, \sporkTidPost, \sporkqtlInit, \sporkQtlAtom, \sporkQtlMask, \sporkAwPost)}
  {}
  {\sporkTidPost \times \{ \sporkqtlInit \} \cup \sporkTid \times \sporkQtlAtom}
  {}
  {\sporkAwPost}

\smallskip
\inOutDiagram[br break=true]
  {\sporkTrArrive(\sporkTidArrive, \sporkQtlArrive, \sporkx, \sporkz_i^*, \sporkz_{\textrm{arrive}})}
  {\sporkTidArrive \times \sporkQtlArrive}
  {}
  {}
  {\{ \sporkPendingAwait(\sporkx, \sporkz_i^*, \sporkz_a) \mid \sporkz_a \ge \sporkz_{\textrm{arrive}}\}}
\hfill
\inOutDiagram
  {\sporkTrAwait(\sporkTidAwait, \sporkQtlAwait, \sporkaw)}
  {}
  {\sporkTidAwait \times \sporkQtlAwait}
  {\{ \sporkaw \}}
  {}

\vspace{3pt} % gaps between diagrams are hand-tuned to look even
\inOutDiagramSides
  {\sporkTrFence(\sporkTidFence, \sporkQtlArrive, \sporkQtlAwait)}
  {\sporkTidFence \times \sporkQtlArrive}
  {\sporkTidFence \times \sporkQtlAwait}
  {}
  {}

\medskip
\vspace{-3pt}
\inOutDiagramSides
  {\sporkTrPreAccess(\sporkmut, \sporkx, \sporkz_i^*, \sporkTidPre, \sporkQtlPre, \sporkAwFree)}
  {\sporkTidPre \times \sporkQtlPre}
  {}
  {\sporkAwFree}
  {}
\caption{Trace event ``inbound sets'' $\sporkObIn$, $\sporkAwIn$ and ``outbound sets'' $\sporkObOut$, $\sporkAwOut$.\\
Note, the use of ``all possible thread IDs'' $\sporkTid$ in $(\sporkTid \times \sporkQtlAtom)$ is intentional.
$\sporkTrPostAccess$ vertices are \emph{sources} with only outbound edges; $\sporkTrPreAccess$ vertices are \emph{sinks} with only inbound edges.
}
\label{fig:InOutSets}
\sporkCodeRule
\end{figure}
